% dynamic.tex -- Section 9: Dynamic Game Creation & Open-Ended Play

\section{Dynamic Game Creation and Open-Ended Play}
\label{sec:dynamic}

Beyond the static game library, Kant supports runtime construction,
registration, and unregistration of games. This enables open-ended play
scenarios where agents themselves can create new strategic environments,
pushing evaluation beyond fixed benchmarks into emergent territory.

% TODO: Motivate open-ended evaluation paradigm

\subsection{Runtime Game Construction API}
\label{sec:game-api}

The dynamic game API provides three construction modes:

\begin{description}[nosep]
    \item[Matrix Construction.] Specify a payoff matrix directly as a nested
        dictionary. The API validates structural consistency (action count
        matching, numeric payoffs) and registers the game.
    \item[Symmetric Construction.] Provide a reduced payoff specification;
        the API expands it into a full symmetric matrix game.
    \item[Custom Payoff Function.] Register a callable that computes payoffs
        from action profiles. Supports arbitrary game logic including
        sequential, stochastic, and information-asymmetric games.
\end{description}

% TODO: API code examples, type signatures

\subsection{Dynamic Registration and Unregistration}
\label{sec:registration}

Games can be registered and unregistered at runtime, allowing the game library
to evolve during a training or evaluation session. The registry validates new
games against schema constraints and assigns unique identifiers. Unregistration
gracefully handles in-progress episodes.

% TODO: Registry protocol, validation schema

\subsection{Open-Ended Play}
\label{sec:open-ended}

In open-ended play mode, agents are given API access to create games mid-session.
This enables several novel evaluation scenarios:

\begin{itemize}[nosep]
    \item \textbf{Game design as capability}: can an agent construct a game
          that incentivizes cooperation?
    \item \textbf{Adversarial game design}: can an agent create games that
          exploit other agents' weaknesses?
    \item \textbf{Negotiated environments}: can agents agree on game rules
          before playing?
    \item \textbf{Curriculum self-play}: can agents create progressively
          more challenging games for self-improvement?
\end{itemize}

% TODO: Experimental protocol for open-ended evaluation

\subsection{Implications for Emergent Behavior and Evaluation}
\label{sec:emergence}

Open-ended play raises fundamental questions about evaluation methodology.
When agents can modify the game space, fixed metrics become insufficient.
We discuss approaches to evaluating emergent behavior, including meta-metrics
that assess the \emph{quality} of games agents create and the stability of
norms that emerge across dynamically generated environments.

% TODO: Meta-evaluation framework, emergent norm analysis
